% Requires: \usepackage{amsmath,amssymb,booktabs,makecell}
% Citation entries: quill_comparison_references.bib
% Insert the table in the main text and the section in the appendix.
% The QUILLS row invokes the paper's main theorem and its numerical conditions.

\begin{table}[t]
\centering
\footnotesize
\setlength{\tabcolsep}{2pt}
\renewcommand{\arraystretch}{1.05}
\caption{Bounds for a fixed univariate analytic class at error $\varepsilon$.
Width is the maximum hidden-layer width; precision covers construction and
evaluation. All bounds are sufficient, not necessary.}
\label{tab:mlp-comparison}
\begin{tabular*}{\linewidth}{@{\extracolsep{\fill}}lcccc@{}}
\toprule
Construction & Width & \makecell{Largest weight\\or bias}
& \makecell{Working precision\\(bits)}
& \makecell{Hidden layers\\/ activation} \\
\midrule
Classical staircase$^{\ddagger}$
& $O(\varepsilon^{-1})$
& $O(\varepsilon^{-1})^{\dagger}$
& $O(\log(1/\varepsilon))^{\dagger}$
& $1$ / tanh \\
Mhaskar~\cite{mhaskar1996}
& $O(\log(1/\varepsilon))$
& \makecell{$\exp\!\big[O(\log^2(1/\varepsilon)$\\
            $\cdot\log\log(1/\varepsilon))\big]^{\dagger}$}
& \makecell{$O(\log^2(1/\varepsilon)$\\
            $\cdot\log\log(1/\varepsilon))^{\dagger}$}
& $1$ / tanh \\
De Ryck et al.~\cite{deryck2021}
& $O(\log(1/\varepsilon))$
& \makecell{$\exp\!\big[O(\log^2(1/\varepsilon)$\\
            $\cdot\log\log(1/\varepsilon))\big]^{\dagger}$}
& \makecell{$O(\log^2(1/\varepsilon)$\\
            $\cdot\log\log(1/\varepsilon))^{\dagger}$}
& $\le 2$ / tanh \\
ChebNet~\cite{tang2019chebnet}
& $O(\log(1/\varepsilon))$
& $O(1)^{\dagger}$
& \makecell{$O(\log(1/\varepsilon)$\\
            $+\log\log(1/\varepsilon))^{\dagger}$}
& \makecell{$O(\log\log(1/\varepsilon))$\\/ ReQU} \\
\textbf{QUILLS}
& $O(\log(1/\varepsilon))$
& $O(\log(1/\varepsilon))$
& \makecell{$O(\log(1/\varepsilon)$\\
            $+\log\log(1/\varepsilon))$}
& $1$ / tanh \\
\bottomrule
\end{tabular*}
\vspace{3pt}
\begin{minipage}{\linewidth}
\scriptsize
$^{\dagger}$Derived for the implementations in
Appendix~\ref{app:comparison-bounds}; not lower bounds.
$^{\ddagger}$Staircase specified there; Jones~\cite{jones1990}
is the historical reference, not the source of these rates.
\end{minipage}
\end{table}

\section{Derivation of the comparison bounds}
\label{app:comparison-bounds}

\paragraph{Scope and arithmetic.}
Fix a uniformly bounded analytic class on a fixed complex neighborhood of
$[-1,1]$. Constants may depend on this class, but not on $\varepsilon$.
We compare normalized $L^2$ error; the uniform bounds used below imply these
$L^2$ bounds. Assume function samples accurate to unit roundoff $u$, accurately
rounded elementary functions, and sufficient exponent range. With $b$
significand bits, $u\asymp 2^{-b}$. If parameter construction and evaluation
together amplify roundoff by at most $A$, then
\begin{equation}
  b \ge \log_2(1/\varepsilon)+\log_2 A+O(1)
  \label{eq:comparison-bits}
\end{equation}
suffices after allocating fixed fractions of the error tolerance to
approximation and arithmetic. This is a sufficient bound, not a necessary
precision requirement. Analyticity gives degree-$n$ Chebyshev approximants
with error $C\rho^{-n}$, for fixed $\rho>1$; hence
$n=O(\log(1/\varepsilon))$.

\paragraph{Classical staircase.}
Set $x_j=-1+2j/N$ and $S(t)=(1+\tanh t)/2$. The network
\begin{equation}
 q_N(x)=f(-1)+\sum_{j=1}^{N}
 [f(x_j)-f(x_{j-1})]S\!\left(N(x-x_j)\right)
\end{equation}
has $N$ neurons. Its coefficient increments are $O(N^{-1})$; the staircase
error and the sum of the exponentially decaying smoothing errors are both
$O(N^{-1})$. Thus $N=O(\varepsilon^{-1})$ suffices. Slopes and biases are
$O(N)$, while total absolute readout weight is bounded. A conservative bound
for construction and evaluation error is $CuN$, so
$b=O(\log(1/\varepsilon))$ suffices.

\paragraph{Mhaskar: an explicit tanh specialization.}
Use the divided-difference approach of~\cite{mhaskar1996}. For $p\ge1$,
approximate $x^p$ by
\begin{equation}
 H_{p,h}(x)=
 \frac{\sum_{r=0}^{p}(-1)^{p-r}\binom{p}{r}
 \tanh\!\left(b_0+(r-p/2)hx\right)}
 {h^p\tanh^{(p)}(b_0)},
 \qquad b_0=\tfrac12\log 2.
\end{equation}
The derivative denominator is nonzero: for $p\ge1$,
$\tanh^{(p)}(b_0)=2^{p+2}A_p(-2)/3^{p+1}$, where
$A_1(z)=1$ and
$A_{p+1}(z)=(1+pz)A_p(z)+z(1-z)A'_p(z)$.
These integer polynomials have constant term $1$, so $A_p(-2)$ is odd and
$|\tanh^{(p)}(b_0)|^{-1}\le\exp(Cp)$.
Taylor's theorem for centered differences gives
\[
 \|H_{p,h}-x^p\|_\infty
 \le \frac{ph^2\|\tanh^{(p+2)}\|_\infty}
 {24|\tanh^{(p)}(b_0)|}.
\]
Cauchy's derivative bound and conversion of the degree-$n$ Chebyshev
approximant to monomials bound the combined error by
$h^2\exp(Cn\log(n+2))$. Choose a dyadic $h$ with
$\log(1/h)=O(\log(1/\varepsilon)+n\log(n+2))$.
All differences share the $2n+1$ slopes $hj/2$, $-n\le j\le n$.
Their combined absolute readout weight is at most $\exp(Cn)(2/h)^n$, giving
\[
 \log B_{\max}
 =O\!\left(\log^2(1/\varepsilon)\log\log(1/\varepsilon)\right).
\]
Stable Chebyshev coefficient computation, monomial conversion, exact integer
derivative numerators, and rounded readout formation incur at most
$u\exp(Cn\log(n+2))(2/h)^n$ error; network evaluation has the same
asymptotic logarithmic amplification. Equation~\eqref{eq:comparison-bits}
therefore gives the stated sufficient bit count.

\paragraph{De Ryck--Lanthaler--Mishra.}
In one dimension, Corollary~5.5 of~\cite{deryck2021} gives widths
$3\lceil s/2\rceil+N-1$ and $6N$, with error at most
$C[3/(2RN)]^s$. Here $R$ is a fixed analyticity radius and $N$ is a partition
size. Fix $N$ so that $3/(2RN)<1$ and take
$s=O(\log(1/\varepsilon))$. The explicit weight estimates in the proof of
Theorem~5.1 give $\log B_{\max}=O(s^2\log(s+2))$.
For construction from samples, first compute a Chebyshev approximant; it is
uniformly bounded on a fixed smaller analytic neighborhood. Its local Taylor
coefficients therefore use degree-independent analyticity constants.
Polynomial basis conversion, the explicit finite-difference scales, and
nonzero activation-derivative denominators give construction amplification
at most $\exp(Cs^2\log(s+2))$.
Two-layer tanh evaluation has error at most
$Cu\,\mathrm{poly}(s)B_{\max}^3$, yielding the stated bit count.
The general construction uses two hidden layers; stronger analyticity
assumptions permit one.

\paragraph{ChebNet.}
The recursive construction in~\cite{tang2019chebnet} represents a degree-$n$
Chebyshev polynomial with $O(n)$ neurons and nonzero parameters and
$O(\log n)$ hidden layers. At each recursive split, a coefficient is doubled
only when its index decreases by at least half. An original coefficient
$c_j$ is therefore amplified by at most $2\max(1,j)$ along any branch.
Since $|c_j|\le C\rho^{-j}$, every intermediate polynomial has coefficient
sum at most $2\sum_j\max(1,j)|c_j|=O(1)$.
The affine/ReQU realization consequently has bounded parameters and
intermediates. Constant fan-in and logarithmic depth give total construction
and evaluation amplification $\mathrm{poly}(n)$, so
$b=\log_2(1/\varepsilon)+O(\log n)$ suffices.

\paragraph{QUILLS.}
This row invokes the main QUILLS approximation and numerical-recovery
theorem, including its bandwidth, aliasing, sampling, and implementation
conditions. It supplies $W=O(\log(1/\varepsilon))$ and controlled readouts.
With grid spacing $h\asymp W^{-1}$, bounded bandwidth $\lambda$, and bounded
center locations, slopes $\gamma=\lambda/h$ and biases are $O(W)$.
The construction-and-evaluation allowance $u\,\mathrm{poly}(W)$ gives
$b=\log_2(1/\varepsilon)+O(\log W)$ via
\eqref{eq:comparison-bits}. The numerical-recovery bound and control of
bandwidth-dependent constants are hypotheses from that theorem; they do not
follow from the width estimate alone.
